% page 468 of the facsimile (image p-467 of the scan)
% block 149.9 x 237.7 mm ; left margin 8.0 mm
\pgc{-12.700mm}{0.000mm}{
\l{\cel{3}460}
\l{}
\l{\sou{prezenta}. I. Qua esas en loko kande ulu arivas an olu, od}
\l{\cel{3}esas en olu. - II. Qua eventas, qua existas en la perio-}
\l{\cel{3}do en qua onu esas. III. (\sou{gram.)Prezento}\cel{1}: Tempo di la}
\l{\cel{3}verbo, qua expresas la tempo-punto ube onu esas. - DEFIS}
\l{}
\l{\sou{prezervar}. (\sou{trans., de, kontre}). Shirmar, protektar, de}
\l{\cel{3}atingeso da malajo. - DEFIRS.}
\l{}
\l{\sou{prezidar}. (\sou{trans.}) Direktar la laboro, la debati da asem-}
\l{\cel{3}blitaro (chambro deliberala, +konselio, e c.) - DEFIRS}
\l{}
\l{\sou{pri}.\cel{2}Koncerne, relate, segun quante koncernesas...}
\l{}
\l{\sou{priklar}. (\sou{netrans.}) Sentar su quaze se pikata neprofunde}
\l{\cel{3}ed iteroze. - DE.}
\l{}
\l{\sou{prima}. Egardata kom origino. - DEFIRS.}
\l{}
\l{\sou{primadono.}\cel{2}Kantistino famoza, tre bone reputata pri lua}
\l{\cel{3}talento. - DEFIRS.}
\l{}
\l{\sou{primara}. I. Qua apartenas a la grado unesma (departante de}
\l{\cel{3}infre). - II.\cel{1}\sou{Kuzi primara}\cel{1}: Qui naskis de patri inter-}
\l{\cel{3}frata. - DEFIS.}
\l{}
\l{\sou{primato}. (\sou{eklezio katol.}) Arkiepiskopo qua dominacas fakte}
\l{\cel{3}ed (od) honore la arkiepiskopi ed episkopi di regiono.-}
\l{\cel{3}DEFIS.}
\l{}
\l{\sou{primino.}\cel{2}(\sou{bot.)}\cel{2}Tegilo unesma di la ovulo (extera).}
\l{}
\l{\sou{primitiva}. Qua aparis ye la origino, e duras havar de ta}
\l{\cel{3}origino, ta o ca karaktero. - DEFIRS.}
\l{}
\l{\sou{primrozo.}\cel{2}(\sou{bot.}) Planto qua florifas en la komenco-dii di}
\l{\cel{3}printempo. - L.\cel{1}\sou{primula grandiflora}. E.}
\l{}
\l{\sou{primulo}. (\sou{bot.}) Narciso sovaja, kun mikra flori. - L.\cel{1}\sou{pri}-}
\l{\cel{3}\sou{mula officinalis}.- DEFIS.}
\l{}
\l{\sou{princo}. I. Persono qua esas suvereno od apartenas a familio}
\l{\cel{3}de suvereni. II. Spozo di princo. - III. (\sou{metaf.}) Per-}
\l{\cel{3}sono qua esas la unesma hierarkiale, la chefo (di la fi-}
\l{\cel{3}lozofi, di la oratori). - DEFIRS.}
\l{}
\l{\sou{principo}. Exaktajo fundamenta, sur qua su apogas la rezono.}
\l{\cel{3}- DEFIRS.}
\l{}
\l{\sou{printempo.}\cel{1}Sezono di la yaro, qua sequas nemediate vintro,}
\l{\cel{3}dum qua la temperaturo di atmosfero divenas plu dolca}
\l{\cel{3}e rinaskas la vejetantaro. (Intertempo qua separas la}
\l{\cel{3}equinoxo qua sequas nemediate vintro, de la solstico di}
\l{\cel{3}somero). - (\sou{anke metaf.})\cel{2}EF.}
}
